% LỜI TỰA --- Quyển XXII
\chapter*{Lời Tựa}
\addcontentsline{toc}{chapter}{Lời Tựa}
\markboth{Lời Tựa}{Lời Tựa}

\begin{truyen}{Lời Tựa}{100 Peak Pedagogy Sentences for Discouraged Students}
\chuhoa{G}{iáo viên bắt đầu bộ sách này với một câu hỏi đơn giản: làm thế nào để mỗi buổi học có thể để lại thứ gì đó đáng nhớ hơn là bài thi?}
\textit{(A teacher began this series with a simple question: how can each class leave behind something more memorable than a test?)}

Khi học sinh nản học, la mắng đôi khi chỉ làm các em đóng cửa thêm~--- nhưng một câu nói đúng: ấm, rõ, không đổ lỗi, có thể mở lại cánh cửa.
\textit{(When students give up on learning, scolding sometimes only makes them shut the door tighter~--- but the right sentence: warm, clear, without blame, can open it again.)}

Quyển hai mươi hai nay được mở rộng thành 100 câu `đỉnh cao sư phạm': nâng đỡ khi học trò mệt, động viên khi các em nghi ngờ mình, chỉ lối khi các em lạc, và giữ các em ở lại với việc học khi lòng đã muốn bỏ.
\textit{(Volume twenty-two is now expanded to 100 `peak pedagogy' sentences: support when students are tired, encouragement when they doubt themselves, direction when they are lost, and words that keep them connected to learning when they feel like giving up.)}

Bộ sách <<Truyện Tích Cảm Hứng Song Ngữ>> gồm 24 quyển, mỗi quyển một chủ đề riêng. Quyển XXII này mang chủ đề: \textbf{nâng đỡ, động viên, mở lối~--- không la mắng}.
\textit{(The series ``Inspiring Bilingual Tales'' contains 24 volumes, each with its own theme. Volume XXII focuses on: \textbf{nâng đỡ, động viên, mở lối~--- không la mắng}.})
\end{truyen}

\begin{baihoc}
Nâng đỡ đúng lúc không làm học trò yếu đuối~--- nó cho các em đủ sức để tự đứng lên.
\textit{(Support at the right time does not make students weak~--- it gives them just enough strength to stand on their own.)}
\end{baihoc}

\begin{ghinhoanh}
\textbf{Short English to remember:}

The right word lifts a student higher than any grade.

Encouragement opens doors that scolding only slams shut.
\end{ghinhoanh}

\begin{tuhocnhanh}
1. Câu nào trong 100 câu này em muốn thầy cô nói với mình nhất lúc em đang nản? \textit{(Which of the 100 sentences would you most want a teacher to say when you are discouraged?)}

2. Em đã từng động viên ai đó học tiếp khi họ muốn bỏ chưa? Bằng lời gì? \textit{(Have you ever encouraged someone to keep studying when they wanted to quit? What did you say?)}

3. Điều gì giúp em tiếp tục khi bản thân không muốn cố gắng nữa? \textit{(What keeps you going when you don't feel like trying anymore?)}
\end{tuhocnhanh}

\begin{center}
\textit{\color{amber}``Nản một lúc không sao; đứng dậy mới quan trọng.''}
\end{center}

\begin{flushright}
\textbf{Tôi Nguyễn Văn Sang}\\
\textit{GV THPT Nguyễn Hữu Cảnh - TP HCM}
\end{flushright}

\ngancach
